\bookchapter{第三十六章：手機同智能家居——令屋企變聰明}

\chapterintro{手機仲可以做一件越嚟越重要嘅事——當成個屋企嘅「遙控器」。由一盞燈、一個插座到一把門鎖，智能家居正在由「極客玩具」變成「家常配置」。呢一章幫你理清協議、生態、中樞同入門套件，再用一份安全清單收尾，令你唔會俾「智能」反噬。}

\booksection{一、智能家居協議：揀邊個好重要}

買智能裝置之前先睇佢支援乜嘢協議，呢個決定咗佢可唔可以俾你嘅手機同生態「差遣」。目前主流協議有四種，各自嘅脾氣如下：

\begin{itemize}
  \item Matter 協議：2022 年由 CSA（連接標準聯盟）發佈，背後係蘋果、Google、Amazon、三星等巨頭聯手。佢嘅目標只有一個——令所有品牌嘅智能裝置講同一種語言。2023-2026 年正快速成為統一標準，支援佢嘅裝置可以同時接入 Apple HomeKit、Google Home、Amazon Alexa 同米家，唔再被單一生態綁架。
  \item Thread 協議：Matter 嘅底層網絡「管道」。低功耗、自組網（某個節點斷電，訊號會自動繞路），反應快、唔卡頓。買有 Thread 邊界路由器嘅裝置（例如 HomePod mini、Apple TV 4K、小米中樞閘道器），體驗明顯更穩。
  \item Wi-Fi：帶寬大、可以直接連路由器，但功耗高、裝置多咗容易擠佔頻道，適合攝像機、掃地機呢類需要傳大數據嘅裝置。
  \item 藍牙（BLE）：配網簡單、慳電，但距離短、唔可以自組網，適合門鎖、溫濕度計呢類小數據裝置。
  \item Zigbee：老牌低功耗自組網協議，要專用閘道器，裝置便宜穩定，但跨品牌互通差，正被 Matter+Thread 逐步替代。
\end{itemize}

\begin{tcbinfo}{一句話原則}
2026 年買智能裝置，優先睇有冇「Matter over Thread」標識——呢個係目前最經得起時間考驗嘅選擇。舊裝置若果只支援 Wi-Fi 或者藍牙都唔使急住淘汰，先用住，新購時再轉去 Matter。

\end{tcbinfo}

\booksection{二、三大生態對比}

協議係「語言」，生態係「圈子」。即使都支援 Matter，你日常仍然要揀一個主生態嚟搭自動化同語音控制。三大主流生態對比如下：

\begin{center}
\begin{tabularx}{\linewidth}{|X|X|X|X|}
\hline
\textbf{生態} & \textbf{優勢} & \textbf{短板} & \textbf{適用人羣} \tabularnewline
\hline
Apple HomeKit & 私隱最強（本地處理為主）、同 iPhone/Siri/快捷指令聯動順滑 & 裝置偏少偏貴、要 HomePod 或者 Apple TV 做中樞 & 全家蘋果裝置、重視私隱同體驗嘅用戶 \tabularnewline
米家（小米） & 裝置最豐富、性價比最高、小愛同學好用 & 生態相對封閉、數據多上雲、廣告偏多 & 預算有限、想一站式配齊全屋嘅入門用戶 \tabularnewline
華為智聯（鴻蒙） & 同鴻蒙手機/平板/手錶無縫流轉、裝置數量快速增長 & 海外裝置支援少、生態仍喺追趕 & 華為手機用戶、睇重「超級終端」多裝置協同 \tabularnewline
\hline
\end{tabularx}
\end{center}

\begin{tcbinfo}{揀生態嘅決策建議}
睇手機揀生態：iPhone 用戶優先 HomeKit，華為用戶優先智聯，其他 Android 用戶米家最慳心。靠 Matter 協議，唔同生態嘅裝置都可以俾同一部手機控制，唔使一次過「企隊」。

\end{tcbinfo}

\booksection{三、手機當家居中樞}

喺多數入門場景入面，手機就係你嘅「家居中樞」——透過 App 控制燈光、冷氣、窗簾、安防，仲可以搭建自動化場景，令屋企「自己鬱起嚟」：

\begin{itemize}
  \item 裝置控制：打開 HomeKit/米家/智慧生活 App，可以遙距或者本地開關燈、調色溫、設定冷氣溫度、拉合窗簾、查看攝像機畫面。
  \item 自動化場景：將「條件動作」組合起嚟，常用嘅三大場景： 返屋企模式：定位到屋企+天黑開玄關燈+開冷氣+拉窗簾+閂門鎖。
  \item 離家模式：最後一個出門（定位離家或者門鎖上鎖）關燈+關冷氣+關插座+啟動攝像機+佈防。
  \item 睡眠模式：每晚 23:00關全屋主燈+拉窗簾+冷氣切 26 度+開夜燈。
\end{itemize}

NFC 碰一碰：喺門貼一張 NFC 標籤，手機一碰就觸發對應場景（碰一下開門燈、再碰一下切睡眠）。iPhone 用「快捷指令」自動化，Android 用「銀包」或者廠商 App 都可以寫入。

\begin{tcbpractice}{嚟試下：搭一個「返屋企模式」}
喺手機智能家居App嘅「自動化」入面新增一條：觸發條件揀「到達某個地點」（屋企），動作依次加「開玄關燈」「冷氣設26度」「閂窗簾」。儲存之後出街行一個圈，返到屋企睇嚇佢會唔會自己著燈——第一次成功嗰陣嘅滿足感，仲過癮過落單買新設備。

\end{tcbpractice}

\booksection{四、入門套件推薦}

唔使一次配齊，由呢四件套開始，就可以覆蓋八成日常需求：

\begin{itemize}
  \item 智能音箱：小愛同學（米家生態，性價比高）/ HomePod mini（蘋果生態，兼做 Thread 中樞）/ Alexa（海外生態廣）。留意：音箱決定咗你嘅語音生態，先定生態再買音箱。
  \item 智能燈泡：Yeelight（米家/Matter，色溫亮度可調，便宜）/ Philips Hue（色彩好、穩定，但要 Hue 閘道器，偏貴）。留意：買之前確認燈座口徑（E27/B22）同係咪需要調色。
  \item 智能插座：小米智能插座（便宜、電量統計）/ TP-Link Kasa（海外 Matter 支援好）。留意：用喺大功率電器（熱水器、電暖器）時務必睇額定功率，唔好過載。
  \item 智能門鎖：鹿客/小米（性價比、米家聯動）/ 凱迪仕（傳統鎖具大廠、機械結構紮實）。留意：確認係咪支援 Matter 或者你所在生態，留意電池續航同應急鎖匙孔位置。
\end{itemize}

\begin{tcbinfo}{選購順序建議}
先買智能音箱（定生態） 再買智能燈泡同插座（低門檻、即刻有反饋） 最後買門鎖同攝像機（涉及安全，謹慎揀）。一次只加一類裝置，熟悉之後再擴，避免「買返嚟唔識用」。

\end{tcbinfo}

\booksection{五、安全注意事項}

智能裝置越多嘅屋企，攻擊面亦越大。下面四條係底線，請逐條對照：

\begin{enumerate}
  \item 改預設密碼：攝像機、路由器、門鎖出廠預設密碼喺網上全部係公開嘅，買到手第一件事就係改成「長而且唔重複」嘅強密碼。
  \item 及時更新固件：廠商推送嘅固件往往修復咗安全漏洞，喺 App 入面開啟「自動更新」，或者至少每月手動檢查一次。
  \item 網絡隔離：將 IoT 裝置單獨放喺一個 Wi-Fi（或者訪客網絡）上，同手機、電腦嘅主網絡分開。萬一某個裝置被攻破，亦難以橫向威脅你嘅重要數據。
  \item 私隱保護：攝像機開啟「私隱區域」遮蔽（例如睡房、洗手間唔影），關閉雲端 24 小時錄像；唔喺屋企嘅攝像機畫面唔好長期開畀第三方 App。
\end{enumerate}

\textbf{呢幾個坑唔好踩}\par

\begin{itemize}
  \item 唔好用門鎖嘅「遙距開鎖」功能畀陌生人——優先用臨時密碼，用完即失效。
  \item 唔好將攝像機直接對住張牀、對住出窗——洩露嘅係你自己嘅私隱。
  \item 唔好貪平買「白牌」無品牌裝置——多數唔更新固件，等於喺屋企裝咗個後門。
\end{itemize}

\begin{tcbsummary}{本章小結}
\begin{itemize}
  \item Matter+Thread 係 2023-2026 年嘅統一方向，買新裝置優先睇呢個標識
  \item 三大生態各有所長：蘋果重私隱體驗、米家重性價比、華為重多裝置協同
  \item 手機可以當中樞，搭好「返屋企/離家/睡眠」三大場景，屋企就可以自己鬱起嚟
  \item 入門四件套：智能音箱、燈泡、插座、門鎖，按呢個順序逐步添置
  \item 改預設密碼、更新固件、網絡隔離、私隱遮蔽——四條安全底線唔可以慳
\end{itemize}

\end{tcbsummary}
